%%%%%%%%%%%%%%%%
%%% exod 23 %%%
%%%%%%%%%%%%%%%%

\Chapter{23}

\Verse{1}
{
	\hE{לֹא תִשָּׂא שֵׁמַע שָׁוְא אַל תָּשֶׁת יָדְךָ עִם רָשָׁע לִהְיֹת עֵד חָמָס׃}
}
{
	\eE{
		23 “You shall not bring up a false report. Do not join\\
		your hand with a wicked person to be a malevolent witness.\\
	}
}
{
}

\Verse{2}
{
	\hE{לֹא תִהְיֶה אַחֲרֵי רַבִּים לְרָעֹת וְלֹא תַעֲנֶה עַל רִב לִנְטֹת אַחֲרֵי רַבִּים לְהַטֹּת׃}
}
{
	\eE{
		You shall not be following many to do bad. And you shall\\
		not testify about a dispute to bend following many, to\\
		bend it.\\
	}
}
{
%	\footnote{
%		not be following many to do bad. Do not follow a group, a crowd, a
%		majority if what they are doing is wrong. Do not do it for acceptance,
%		for the secure feeling of being in a group, or for the sadistic pleasure
%		of being able to exclude someone. It is easy to be hurtful in a group.
%		And it is easy to keep silent when one’s group does harm—or when its
%		leaders do harm from their position, which derives its power from the
%		group. All of this is forbidden. It is utterly inconsistent with the
%		Torah’s conceptions of what a human should be and how one should behave
%		toward other human beings. This comment should be unnecessary. Footnote
%		23:2. to bend following many, to bend it. When giving testimony or when
%		sitting in judgment, do not incline toward the majority of those
%		involved in the case or toward the majority of the public. To “bend”
%		toward them is to “bend” justice.
%	}
}

\Verse{3}
{
	\hE{וְדָל לֹא תֶהְדַּר בְּרִיבוֹ׃}
}
{
	\eE{And you shall not favor a weak person in his dispute.}
}
{
%	\footnote{
%		you shall not favor a weak person. Of course one should not favor the
%		majority, or the rich, or the famous, or the powerful. Of course we
%		should not bend the law for bad motives. We are reminded that we also
%		must not bend the law or the truth even for good motives, such as
%		feelings of compassion for the poor. Compare Lev 19:15.
%	}
}

\Verse{4}
{
	\hE{כִּי תִפְגַּע שׁוֹר אֹיִבְךָ אוֹ חֲמֹרוֹ תֹּעֶה הָשֵׁב תְּשִׁיבֶנּוּ לוֹ׃}
}
{
	\eE{
		“If you will happen upon your enemy’s ox or his ass\\
		straying, you shall bring it back to him.\\
	}
}
{
}

\Verse{5}
{
	\hE{כִּי תִרְאֶה חֲמוֹר שֹׂנַאֲךָ רֹבֵץ תַּחַת מַשָּׂאוֹ וְחָדַלְתָּ מֵעֲזֹב לוֹ עָזֹב תַּעֲזֹב עִמּוֹ׃}
}
{
	\eE{
		If you will see the ass of someone who hates you sagging\\
		under its burden, and you would hold back from helping\\
		him: you shall help with him.\\
	}
}
{
%	\footnote{
%		you would hold back from helping him. That is, you would hold back from
%		helping your adversary. Footnote 23:5. help. The Hebrew ‘zb (to leave)
%		is difficult. One possibility, suggested here, is to read the word as
%		‘zr (to help). Alternatively, it may mean: “If you will see the ass of
%		someone who hates you sagging under its burden, then you shall hold back
%		from leaving it for him; you shall leave it with him.” That is, do not
%		leave it for him to handle himself; do what is necessary to assist and
%		leave it back with him. By any reading, though, the main point is that
%		one must be of help, even to someone who bears one ill will—perhaps
%		especially to someone who bears ill will.
%	}
}

\Verse{6}
{
	\hE{לֹא תַטֶּה מִשְׁפַּט אֶבְיֹנְךָ בְּרִיבוֹ׃}
}
{
	\eE{
		“You shall not bend the judgment of your poor in his\\
		dispute.\\
	}
}
{
%	\footnote{
%		your poor. It is natural to use the phrase “your enemy,” as in the
%		preceding verses, but we would not readily expect “your poor”
%		linguistically. But the issue is not linguistic. It is moral. The poor
%		are not neutral as a group. They are your poor, your responsibility.
%	}
}

\Verse{7}
{
	\hE{מִדְּבַר שֶׁקֶר תִּרְחָק וְנָקִי וְצַדִּיק אַל תַּהֲרֹג כִּי לֹא אַצְדִּיק רָשָׁע׃}
}
{
	\eE{
		“You shall keep far from a word of a lie, and do not kill\\
		an innocent and a virtuous person, because I shall not\\
		vindicate a wicked person.\\
	}
}
{
}

\Verse{8}
{
	\hE{וְשֹׁחַד לֹא תִקָּח כִּי הַשֹּׁחַד יְעַוֵּר פִּקְחִים וִיסַלֵּף דִּבְרֵי צַדִּיקִים׃}
}
{
	\eE{
		“And you shall not take a bribe, because bribery will\\
		blind those who can see and will undermine the words of\\
		the virtuous.\\
	}
}
{
}

\Verse{9}
{
	\hE{וְגֵר לֹא תִלְחָץ וְאַתֶּם יְדַעְתֶּם אֶת נֶפֶשׁ הַגֵּר כִּי גֵרִים הֱיִיתֶם בְּאֶרֶץ מִצְרָיִם׃}
}
{
	\eE{
		“And you shall not oppress an alien—since you know the\\
		alien’s soul, because you were aliens in the land of\\
		Egypt.\\
	}
}
{
%	\footnote{
%		you know the alien’s soul. This verse, commanding not to oppress an
%		alien residing among Israel, repeats a commandment that has already been
%		given just a little earlier, in 22:20. But it says it even more
%		eloquently. You have been aliens, so you know the alien. You know how he
%		or she feels oppression inside. Few of the laws in this code have an
%		extra line of explanation like this. Notably, another law protecting
%		aliens in the preceding chapter also adds a line of explanation (22:20).
%		Again, the treatment of a non-Israelite who is within Israel’s power is
%		singled out in a special way.
%	}
}

\Verse{10}
{
	\hE{וְשֵׁשׁ שָׁנִים תִּזְרַע אֶת אַרְצֶךָ וְאָסַפְתָּ אֶת תְּבוּאָתָהּ׃}
}
{
	\eE{
		“And six years you shall sow your land and gather its\\
		produce;\\
	}
}
{
}

\Verse{11}
{
	\hE{וְהַשְּׁבִיעִת תִּשְׁמְטֶנָּה וּנְטַשְׁתָּהּ וְאָכְלוּ אֶבְיֹנֵי עַמֶּךָ וְיִתְרָם תֹּאכַל חַיַּת הַשָּׂדֶה כֵּן תַּעֲשֶׂה לְכַרְמְךָ לְזֵיתֶךָ׃}
}
{
	\eE{
		and the seventh: you shall let it lie fallow and leave it,\\
		and your people’s indigent will eat, and what they leave\\
		the animal of the field will eat. You shall do this to\\
		your vineyard, to your olives.\\
	}
}
{
}

\Verse{12}
{
	\hE{שֵׁשֶׁת יָמִים תַּעֲשֶׂה מַעֲשֶׂיךָ וּבַיּוֹם הַשְּׁבִיעִי תִּשְׁבֹּת לְמַעַן יָנוּחַ שׁוֹרְךָ וַחֲמֹרֶךָ וְיִנָּפֵשׁ בֶּן אֲמָתְךָ וְהַגֵּר׃}
}
{
	\eE{
		“Six days you shall do the things you do, and on the\\
		seventh day you shall cease, so that your ox and your ass\\
		will rest, and your maid’s son and the alien will be\\
		refreshed.\\
	}
}
{
}

\Verse{13}
{
	\hE{וּבְכֹל אֲשֶׁר אָמַרְתִּי אֲלֵיכֶם תִּשָּׁמֵרוּ וְשֵׁם אֱלֹהִים אֲחֵרִים לֹא תַזְכִּירוּ לֹא יִשָּׁמַע עַל פִּיךָ׃}
}
{
	\eE{
		“And you shall be watchful in everything that I have said\\
		to you. And you shall not commemorate the name of other\\
		gods. Let it not be heard on your mouth.\\
	}
}
{
%	\footnote{
%		commemorate. Other translations take this to mean “mention,” i.e., that
%		one cannot even say the name of another god—but compare Exod 20:21,
%		where this same word appears to refer to mentioning YHWH’s name
%		institutionally. I therefore understand this verse to mean that one is
%		commanded not to say another god’s name in a setting of worship.
%	}
}

\Verse{14}
{
	\hE{שָׁלֹשׁ רְגָלִים תָּחֹג לִי בַּשָּׁנָה׃}
}
{
	\eE{“You shall celebrate three pilgrimages for me in the year:}
}
{
}

\Verse{15}
{
	\hE{אֶת חַג הַמַּצּוֹת תִּשְׁמֹר שִׁבְעַת יָמִים תֹּאכַל מַצּוֹת כַּאֲשֶׁר צִוִּיתִךָ לְמוֹעֵד חֹדֶשׁ הָאָבִיב כִּי בוֹ יָצָאתָ מִמִּצְרָיִם וְלֹא יֵרָאוּ פָנַי רֵיקָם׃}
}
{
	\eE{
		“You shall observe the Festival of Unleavened Bread. Seven\\
		days you shall eat unleavened bread, as I commanded you,\\
		at the appointed time of the month of Abib, because you\\
		went out of Egypt in it. And none shall appear before me\\
		empty.\\
	}
}
{
%	\footnote{
%		empty. Meaning: empty-handed. One must go to the designated place and
%		bring a sacrifice.
%	}
}

\Verse{16}
{
	\hE{וְחַג הַקָּצִיר בִּכּוּרֵי מַעֲשֶׂיךָ אֲשֶׁר תִּזְרַע בַּשָּׂדֶה וְחַג הָאָסִף בְּצֵאת הַשָּׁנָה בְּאסְפְּךָ אֶת מַעֲשֶׂיךָ מִן הַשָּׂדֶה׃}
}
{
	\eE{
		“And the Festival of Harvest, the first fruits of what you\\
		do, of what you will sow in the field. “And the Festival\\
		of Gathering, at the end of the year, when you gather what\\
		you have done from the field.\\
	}
}
{
}

\Verse{17}
{
	\hE{שָׁלֹשׁ פְּעָמִים בַּשָּׁנָה יֵרָאֶה כּל זְכוּרְךָ אֶל פְּנֵי הָאָדֹן יְהֹוָה׃}
}
{
	\eE{
		“Three times in the year every male of yours shall appear\\
		in front of the Lord YHWH.\\
	}
}
{
}

\Verse{18}
{
	\hE{לֹא תִזְבַּח עַל חָמֵץ דַּם זִבְחִי וְלֹא יָלִין חֵלֶב חַגִּי עַד בֹּקֶר׃}
}
{
	\eE{
		“You shall not offer the blood of my sacrifice on leavened\\
		bread. And the fat of my festival shall not remain until\\
		morning.\\
	}
}
{
}

\Verse{19}
{
	\hE{רֵאשִׁית בִּכּוּרֵי אַדְמָתְךָ תָּבִיא בֵּית יְהֹוָה אֱלֹהֶיךָ לֹא תְבַשֵּׁל גְּדִי בַּחֲלֵב אִמּוֹ׃}
}
{
	\eE{
		“You shall bring the first of the firstfruits of your land\\
		to the house of YHWH, your God. “You shall not cook a kid\\
		in its mother’s milk.\\
	}
}
{
%	\footnote{
%		the house of YHWH. This is a reference to the Temple, the building that
%		will one day be the only sanctioned place of sacrifice and the central
%		place of worship. The first house of YHWH will be located at Shiloh. The
%		second and third will be located at Jerusalem. Footnote 23:19. a kid in
%		its mother’s milk. A Canaanite (Ugaritic) text may picture the chief
%		god, El, having kid cooked in milk. This biblical law may thus be a
%		prohibition against eating something that was regarded as a food for a
%		deity. This explanation is uncertain because the Canaanite text does not
%		specify mother’s milk, and the word there that is read as “kid” is
%		partly effaced. Alternatively, the biblical law may be a moral one,
%		consistent with other biblical law, which prohibits sacrificing an ox or
%		sheep on the same day as its offspring (Lev 22:28) or taking a mother
%		bird along with her eggs in the wild (Deut 22:6–7). This law came to be
%		taken in later Jewish practice to amount to a complete separation of
%		meat and dairy consumption. This practice is not commanded in the Torah
%		but is rather based on postbiblical rabbinic rulings.
%	}
}

\Verse{20}
{
	\hJ{הִנֵּה אָנֹכִי שֹׁלֵחַ מַלְאָךְ לְפָנֶיךָ לִשְׁמרְךָ בַּדָּרֶךְ וְלַהֲבִיאֲךָ אֶל הַמָּקוֹם אֲשֶׁר הֲכִנֹתִי׃}
}
{
	\eJ{
		“Here, I’m sending an angel ahead of you to watch over you\\
		on the way and to bring you to the place that I’ve\\
		prepared.\\
	}
}
{
%	\footnote{
%		I’m sending an angel ahead of you. This will be repeated twice (32:34;
%		33:2). It is another recurrence of language from the accounts of the
%		patriarchs: Abraham told his servant that God will send “His angel ahead
%		of you,” which would lead to the success of his mission to get a wife
%		for Isaac (Gen 24:7; see also 32:4). That is now a reminder that the
%		preceding angel on a journey is an assurance and protection.
%	}
}

\Verse{21}
{
	\hJ{הִשָּׁמֶר מִפָּנָיו וּשְׁמַע בְּקֹלוֹ אַל תַּמֵּר בּוֹ כִּי לֹא יִשָּׂא לְפִשְׁעֲכֶם כִּי שְׁמִי בְּקִרְבּוֹ׃}
}
{
	\eJ{
		Be watchful in front of him and listen to his voice. Don’t\\
		rebel against him, because he will not bear your offense,\\
		because my name is within him.\\
	}
}
{
%	\footnote{
%		rebel. The Hebrew, tammr, here in the Masoretic Text has the form of a
%		Hiphil of the root mrr, meaning “to make bitter.” But these vowel points
%		are generally understood to be incorrect. The word should rather be read
%		as temer, a form of the root mrh, meaning “to rebel,” which is how the
%		Greek reads. The theme of rebellion will figure strongly in the events
%		to come. Alternatively, one might argue that it is an allusion to the
%		story of the bitter waters at Marah, which was the people’s first
%		rebellion after crossing the Red Sea.
%	}
}

\Verse{22}
{
	\hJ{כִּי אִם שָׁמוֹעַ תִּשְׁמַע בְּקֹלוֹ וְעָשִׂיתָ כֹּל אֲשֶׁר אֲדַבֵּר וְאָיַבְתִּי אֶת אֹיְבֶיךָ וְצַרְתִּי אֶת צֹרְרֶיךָ׃}
}
{
	\eJ{
		But if you will listen to his voice and do everything that\\
		I speak, then I’ll be an enemy to your enemies and an\\
		opponent to your opponents.\\
	}
}
{
}

\Verse{23}
{
	\hJ{כִּי יֵלֵךְ מַלְאָכִי לְפָנֶיךָ וֶהֱבִיאֲךָ אֶל הָאֱמֹרִי וְהַחִתִּי וְהַפְּרִזִּי וְהַכְּנַעֲנִי הַחִוִּי וְהַיְבוּסִי וְהִכְחַדְתִּיו׃}
}
{
	\eJ{
		When my angel will go ahead of you and will bring you to\\
		the Amorite and the Hittite and the Perizzite and the\\
		Canaanite, the Hivite, and the Jebusite, and I obliterate\\
		them,\\
	}
}
{
}

\Verse{24}
{
	\hJ{לֹא תִשְׁתַּחֲוֶה לֵאלֹהֵיהֶם וְלֹא תעבְדֵם וְלֹא תַעֲשֶׂה כְּמַעֲשֵׂיהֶם כִּי הָרֵס תְּהָרְסֵם וְשַׁבֵּר תְּשַׁבֵּר מַצֵּבֹתֵיהֶם׃}
}
{
	\eJ{
		you shall not bow to their gods, and you shall not serve\\
		them, and you shall not do like the things they do; but\\
		you shall tear them down and shatter their pillars.\\
	}
}
{
}

\Verse{25}
{
	\hJ{וַעֲבַדְתֶּם אֵת יְהֹוָה אֱלֹהֵיכֶם וּבֵרַךְ אֶת לַחְמְךָ וְאֶת מֵימֶיךָ וַהֲסִרֹתִי מַחֲלָה מִקִּרְבֶּךָ׃}
}
{
	\eJ{
		And you will serve YHWH, your God, and He will bless your\\
		bread and your water, and I shall turn sickness away from\\
		within you.\\
	}
}
{
}

\Verse{26}
{
	\hJ{לֹא תִהְיֶה מְשַׁכֵּלָה וַעֲקָרָה בְּאַרְצֶךָ אֶת מִסְפַּר יָמֶיךָ אֲמַלֵּא׃}
}
{
	\eJ{
		There won’t be a bereaved woman and an infertile woman in\\
		your land. I shall fulfill the number of your days.\\
	}
}
{
}

\Verse{27}
{
	\hJ{אֶת אֵימָתִי אֲשַׁלַּח לְפָנֶיךָ וְהַמֹּתִי אֶת כּל הָעָם אֲשֶׁר תָּבֹא בָּהֶם וְנָתַתִּי אֶת כּל אֹיְבֶיךָ אֵלֶיךָ עֹרֶף׃}
}
{
	\eJ{
		I shall have my terror go ahead of you, and I shall throw\\
		all the people against whom you come into tumult, and I\\
		shall give all your enemies to you by the back.\\
	}
}
{
%	\footnote{
%		by the back. The enemies will turn their backs and flee. This image of
%		Israel’s enemies occurs in 2 Sam 22:41 and Ps 18:41. As a terrifying
%		lesson, it is reversed in a case of a serious violation by an Israelite
%		of a divine command, and then it is Israel that turns its back to its
%		enemies (Josh 7:8,12).
%	}
}

\Verse{28}
{
	\hJ{וְשָׁלַחְתִּי אֶת הַצִּרְעָה לְפָנֶיךָ וְגֵרְשָׁה אֶת הַחִוִּי אֶת הַכְּנַעֲנִי וְאֶת הַחִתִּי מִלְּפָנֶיךָ׃}
}
{
	\eJ{
		And I shall send the hornet ahead of you, and it will\\
		drive out the Hivite, the Canaanite, and the Hittite from\\
		in front of you.\\
	}
}
{
}

\Verse{29}
{
	\hJ{לֹא אֲגָרְשֶׁנּוּ מִפָּנֶיךָ בְּשָׁנָה אֶחָת פֶּן תִּהְיֶה הָאָרֶץ שְׁמָמָה וְרַבָּה עָלֶיךָ חַיַּת הַשָּׂדֶה׃}
}
{
	\eJ{
		I won’t drive them out from in front of you in one year,\\
		or else the land will be a devastation, and the animal of\\
		the field will be many at you.\\
	}
}
{
}

\Verse{30}
{
	\hJ{מְעַט מְעַט אֲגָרְשֶׁנּוּ מִפָּנֶיךָ עַד אֲשֶׁר תִּפְרֶה וְנָחַלְתָּ אֶת הָאָרֶץ׃}
}
{
	\eJ{
		Little by little I shall drive them out from in front of\\
		you, until you will be fruitful, and you will have a\\
		legacy of the land.\\
	}
}
{
}

\Verse{31}
{
	\hJ{וְשַׁתִּי אֶת גְּבֻלְךָ מִיַּם סוּף וְעַד יָם פְּלִשְׁתִּים וּמִמִּדְבָּר עַד הַנָּהָר כִּי אֶתֵּן בְּיֶדְכֶם אֵת יֹשְׁבֵי הָאָרֶץ וְגֵרַשְׁתָּמוֹ מִפָּנֶיךָ׃}
}
{
	\eJ{
		And I shall set your border from the Red Sea to the sea of\\
		the Philistines and from the wilderness to the river,\\
		because I shall put the residents of the land in your\\
		hand, and you’ll drive them out from in front of you.\\
	}
}
{
%	\footnote{
%		the Red Sea. This refers to the eastern arm of the Red Sea, presently
%		the site of Eilat. The promised territory of Israel is never conceived
%		of in the Torah as extending all the way across the Sinai Peninsula to
%		the western arm of the Red Sea. It stops at the “river of Egypt,” the
%		Wadi el-‘Arîsh, which is north of the eastern arm of the Red Sea.
%		Footnote 23:31. the sea of the Philistines. The Mediterranean Sea. The
%		Philistines came across the Mediterranean from the Greek islands. They
%		lived in cities near the Mediterranean coast (Ashdod, Ashkelon, Ekron,
%		Gath, and Gaza). Footnote 23:31. the river. This is commonly understood
%		to refer to the Euphrates.
%	}
}

\Verse{32}
{
	\hJ{לֹא תִכְרֹת לָהֶם וְלֵאלֹהֵיהֶם בְּרִית׃}
}
{
	\eJ{
		You shall not make a covenant with them and with their\\
		gods.\\
	}
}
{
}

\Verse{33}
{
	\hJ{לֹא יֵשְׁבוּ בְּאַרְצְךָ פֶּן יַחֲטִיאוּ אֹתְךָ לִי כִּי תַעֲבֹד אֶת אֱלֹהֵיהֶם כִּי יִהְיֶה לְךָ לְמוֹקֵשׁ׃}
}
{
	\eJ{
		They shall not live in your land, in case they will make\\
		you sin against me when you’ll serve their gods, because\\
		it will be a trap for you.”\\
	}
}
{
}
